%% ============================================================
%%  Artigo Científico — Disciplina de Pesquisa e Ordenação
%%  Universidade Tecnológica Federal do Paraná — UTFPR
%%  Formato: SBC (Sociedade Brasileira de Computação)
%% ============================================================

\documentclass[12pt]{article}

\usepackage{sbc-template}
\usepackage{graphicx,url}

% Diretório base para imagens — funciona independente de onde o
% compilador é invocado (dentro de artigo/ ou da raiz do projeto)
\graphicspath{
  {../output/plots_python/}%  quando compilado de dentro de artigo/
  {output/plots_python/}%     quando compilado da raiz do projeto
}
% babel com português — use texlive-lang-portuguese no seu sistema
% \usepackage[brazil]{babel}
\usepackage[nil,shorthands=off]{babel}
\usepackage[utf8]{inputenc}
\usepackage{booktabs}       % tabelas com melhor tipografia
\usepackage{amsmath}        % notação matemática
\usepackage{multirow}       % células mescladas em tabelas

% Substitutos para siunitx
\newcommand{\num}[1]{#1}

\sloppy

% ──────────────────────────────────────────────────────────────
%  Metadados
% ──────────────────────────────────────────────────────────────
\title{Análise Empírica de Algoritmos de Ordenação:\\
       Uma Avaliação Experimental de Bubble Sort, Selection Sort,\\
       Insertion Sort e Shell Sort em Entradas de Grande Volume}

\author{Gustavo R. Mazur\inst{1}}

\address{
  Departamento de Informática -- Universidade Tecnológica Federal do Paraná (UTFPR)\\
  \email{gustavomazur@alunos.utfpr.edu.br}
}

\begin{document}

\maketitle

% ──────────────────────────────────────────────────────────────
%  Resumo e Abstract
% ──────────────────────────────────────────────────────────────
\begin{abstract}
Sorting algorithms are among the most fundamental structures in
Computer Science, with well-established theoretical complexities.
However, asymptotic analysis abstracts hardware constants, memory
effects, and input distribution, which can cause significant
divergences between theoretical predictions and observed behavior.
This work presents an empirical study comparing four classical
sorting algorithms — Bubble Sort, Selection Sort, Insertion Sort,
and Shell Sort — through a reproducible benchmarking infrastructure
implemented in Go~1.23. Experiments were conducted on five input
sizes (175,000 to 1,500,000 elements) under three dataset
configurations (sorted, reversed, and random), with three
independent runs per combination. Metrics collected include
execution time, elementary operations (comparisons and swaps),
and heap memory allocations. The results empirically confirm
theoretical complexity predictions and reveal relevant practical
distinctions: Bubble Sort and Insertion Sort achieve near-zero
execution time for sorted inputs; Selection Sort maintains
constant behavior regardless of input order; and Shell Sort
outperforms all $O(n^2)$ algorithms by up to \num{12657} times
for random inputs of 750,000 elements, demonstrating that
complexity class membership is the dominant factor in
large-scale performance.

\noindent{\bf Keywords:} Sorting Algorithms, Empirical Analysis, Computational Complexity, Go Language.
\end{abstract}

\begin{resumo}
Algoritmos de ordenação figuram entre as estruturas mais
fundamentais da Ciência da Computação, com complexidades teóricas
amplamente estabelecidas. Contudo, a análise assintótica abstrai
constantes de hardware, efeitos de memória e distribuição de
entrada, o que pode gerar divergências significativas entre
predições teóricas e comportamento observado. O presente trabalho
apresenta um estudo empírico comparando quatro algoritmos de
ordenação clássicos — Bubble Sort, Selection Sort, Insertion Sort
e Shell Sort — por meio de uma infraestrutura de benchmarking
reprodutível implementada em Go~1.23. Os experimentos foram
conduzidos sobre cinco tamanhos de entrada (175.000 a 1.500.000
elementos) sob três configurações de dataset (ordenado, invertido
e aleatório), com três execuções independentes por combinação.
As métricas coletadas incluem tempo de execução, operações
elementares (comparações e trocas) e alocações de memória em
\textit{heap}. Os resultados confirmam empiricamente as predições
de complexidade teórica e revelam distinções práticas relevantes:
Bubble Sort e Insertion Sort atingem tempo de execução
praticamente nulo para entradas ordenadas; Selection Sort mantém
comportamento constante independentemente da ordem da entrada; e
Shell Sort supera os demais algoritmos em até \num{12657} vezes
para entradas aleatórias com 750.000 elementos, demonstrando que
a classe de complexidade é o fator dominante no desempenho
em grande escala.

\noindent{\bf Palavras-chave:} Algoritmos de Ordenação, Análise Empírica, Complexidade Computacional, Linguagem Go.
\end{resumo}

% ──────────────────────────────────────────────────────────────
%  1. INTRODUÇÃO
% ──────────────────────────────────────────────────────────────
\section{Introdução}

A ordenação de dados constitui um dos problemas centrais da
Ciência da Computação. Algoritmos de ordenação estão presentes
em praticamente todas as camadas de sistemas de software, desde
rotinas de banco de dados até renderização gráfica e indexação
de mecanismos de busca \cite{cormen2009}. A importância desse
problema motivou décadas de pesquisa, resultando em algoritmos
com complexidades assintóticas bem conhecidas e provadas.

A análise de complexidade assintótica, expressa em notação
$O(\cdot)$, descreve o comportamento de um algoritmo no limite,
quando o tamanho da entrada $n$ tende ao infinito. Essa
abstração é teoricamente rigorosa, porém deliberadamente ignora
constantes multiplicativas, coeficientes de ordem inferior e
efeitos de hardware, tais como localidade de \textit{cache},
\textit{pipeline} de processador e custo de alocação de memória
\cite{knuth1998}. Em contextos práticos, esses fatores podem
produzir diferenças de desempenho substanciais entre algoritmos
pertencentes à mesma classe de complexidade.

Diante desse cenário, o presente trabalho adota uma abordagem
complementar à análise teórica: a avaliação empírica e
reprodutível do desempenho de algoritmos de ordenação em
condições reais de hardware. Especificamente, este estudo tem
os seguintes objetivos:

\begin{enumerate}
  \item Implementar e instrumentar quatro algoritmos clássicos
        de ordenação com contagem exata de comparações e trocas;
  \item Executar os algoritmos sobre três tipos de dataset
        (ordenado, invertido e aleatório) em cinco tamanhos de
        entrada (175.000 a 1.500.000 elementos);
  \item Comparar os resultados empíricos com as predições
        teóricas de complexidade;
  \item Identificar e discutir divergências e nuances práticas
        não capturadas pela análise assintótica.
\end{enumerate}

Os algoritmos avaliados são: Bubble Sort, Selection Sort,
Insertion Sort e Shell Sort. Esses quatro algoritmos foram
selecionados por representarem uma progressão de complexidade
e estratégias distintas de ordenação, permitindo uma análise
comparativa abrangente dentro da família de algoritmos
baseados em comparação.

O restante deste artigo está organizado da seguinte forma:
a Seção~\ref{sec:theory} apresenta o referencial teórico dos
algoritmos avaliados; a Seção~\ref{sec:related} discute os
trabalhos relacionados; a Seção~\ref{sec:method} descreve a
metodologia experimental adotada; a Seção~\ref{sec:results}
apresenta os resultados obtidos; a Seção~\ref{sec:discussion}
analisa e discute esses resultados; e a Seção~\ref{sec:conclusion}
conclui o artigo.

% ──────────────────────────────────────────────────────────────
%  2. REFERENCIAL TEÓRICO
% ──────────────────────────────────────────────────────────────
\section{Referencial Teórico}
\label{sec:theory}

\subsection{Bubble Sort}

O Bubble Sort é um dos algoritmos de ordenação mais simples e
didaticamente reconhecidos. Seu funcionamento baseia-se em
passagens sucessivas sobre o vetor, comparando e trocando
pares de elementos adjacentes quando encontrados fora de ordem
\cite{cormen2009}. A cada passagem completa, o maior elemento
não posicionado é propagado ("borbulhado") ao final do vetor.

A complexidade de tempo do Bubble Sort é $O(n^2)$ no caso médio
e no pior caso, e $O(n)$ no melhor caso, quando o vetor já se
encontra ordenado — condição detectável por uma \textit{flag}
de troca \cite{cormen2009}. No pior caso, para um vetor de $n$
elementos em ordem invertida, a implementação não otimizada
realiza $n(n-1)$ comparações, pois cada passagem percorre
o vetor inteiro sem redução do limite superior. A complexidade
de espaço é $O(1)$, pois o algoritmo opera \textit{in-place}.

\subsection{Selection Sort}

O Selection Sort divide conceitualmente o vetor em duas partes:
uma porção ordenada (à esquerda) e uma não ordenada (à direita).
A cada iteração, o algoritmo seleciona o menor elemento da
porção não ordenada e o posiciona ao final da porção ordenada,
mediante uma única troca \cite{cormen2009}.

Uma característica fundamental do Selection Sort é a ausência
de melhor caso em termos de comparações: independentemente
da distribuição de entrada, o algoritmo realiza exatamente
$\frac{n(n-1)}{2}$ comparações, o que implica complexidade
$O(n^2)$ para todos os casos. Em contrapartida, o número
máximo de trocas é $n-1$ — propriedade vantajosa em cenários
onde a operação de escrita é significativamente mais custosa
que a de leitura \cite{knuth1998}. A complexidade de espaço
é $O(1)$.

\subsection{Insertion Sort}

O Insertion Sort constrói o vetor ordenado incrementalmente:
a cada iteração, o algoritmo insere o elemento $i$-ésimo na
posição correta dentro dos $i-1$ elementos já ordenados,
deslocando os demais elementos à direita \cite{cormen2009}.

A complexidade de tempo é $O(n^2)$ no pior caso (vetor invertido)
e $O(n)$ no melhor caso (vetor já ordenado), pois neste último
cada elemento requer exatamente uma comparação para confirmar
sua posição. O Insertion Sort é particularmente eficiente para
entradas pequenas ou parcialmente ordenadas, sendo o algoritmo
de escolha para finalização em implementações híbridas modernas,
como o Tim Sort, utilizado em Python e Java \cite{cormen2009}.
A complexidade de espaço é $O(1)$.

\subsection{Shell Sort}

O Shell Sort, proposto por Donald Shell em 1959, estende o
Insertion Sort introduzindo o conceito de \textit{gap}: em vez
de comparar apenas elementos adjacentes, o algoritmo compara e
ordena elementos separados por um intervalo $h$, reduzindo
progressivamente esse intervalo até $h = 1$ \cite{shell1959}.
Essa estratégia permite mover elementos distantes de sua posição
correta em poucas operações, reduzindo o número total de trocas.

A complexidade do Shell Sort depende da sequência de gaps
utilizada. Com a sequência original de Shell ($h = \lfloor n/2
\rfloor, \lfloor n/4 \rfloor, \ldots, 1$), a complexidade de
pior caso é $O(n^2)$, mas o comportamento prático é
significativamente melhor, observado empiricamente como
$O(n \log^2 n)$ em média \cite{knuth1998}. Sequências de gaps
mais sofisticadas, como as de Hibbard e Sedgewick, garantem
complexidade $O(n^{3/2})$ e $O(n^{4/3})$ no pior caso,
respectivamente \cite{sedgewick1996}. A complexidade de espaço
é $O(1)$.

% ──────────────────────────────────────────────────────────────
%  3. TRABALHOS RELACIONADOS
% ──────────────────────────────────────────────────────────────
\section{Trabalhos Relacionados}
\label{sec:related}

A literatura sobre análise experimental de algoritmos de
ordenação pode ser organizada em dois grupos principais:
estudos empíricos comparativos de desempenho e análises
específicas de algoritmos e sua interação com características
de hardware.

\subsection*{Estudos Empíricos Comparativos de Desempenho}

Pizarro-Vasquez et al.~\cite{pizarro2020} conduziram uma
avaliação empírica de onze algoritmos de ordenação —
incluindo Bubble Sort, Insertion Sort, Selection Sort,
Shell Sort, Merge Sort, Heap Sort e Quick Sort —
implementados em C e Python. Os experimentos utilizaram
quatro configurações de entrada (aleatória, quase ordenada,
invertida e com poucos valores únicos), avaliando tempos
de execução. Os autores confirmaram que algoritmos de classe
$O(n \log n)$ superam consistentemente os de classe $O(n^2)$
à medida que o tamanho da entrada cresce. Contudo, o estudo
não instrumenta contagens exatas de operações elementares
nem quantifica a variabilidade estatística entre execuções
independentes, aspectos que o presente trabalho aborda
explicitamente.

Biggar et al.~\cite{biggar2008} investigaram o comportamento
de algoritmos de ordenação clássicos sob a perspectiva de
predição de desvios (\textit{branch prediction}) em
processadores modernos. Os autores demonstraram que algoritmos
com mesma complexidade assintótica podem divergir
substancialmente em desempenho prático em função de
características microarquiteturais, como \textit{miss rate}
do preditor de desvios e localidade de \textit{cache}. Esse
estudo evidencia que fatores de hardware — abstraídos pela
notação $O(\cdot)$ — são determinantes no desempenho
observado em condições reais, hipótese que também motiva a
abordagem experimental adotada neste trabalho.

\subsection*{Análise Teórico-Empírica do Shell Sort}

Sedgewick~\cite{sedgewick1996} apresentou uma análise
abrangente do Shell Sort e das sequências de gaps que
determinam sua complexidade. O trabalho demonstrou que a
escolha da sequência é determinante: a sequência original
de Shell resulta em $O(n^2)$ no pior caso, enquanto as
sequências de Hibbard e Sedgewick garantem $O(n^{3/2})$
e $O(n^{4/3})$, respectivamente. A análise fornece
embasamento teórico para a interpretação do comportamento
empírico do Shell Sort investigado neste artigo.

Diferentemente dos trabalhos acima, o presente artigo
adota instrumentação \textit{in-code} que contabiliza
exatamente as comparações e trocas realizadas por cada
algoritmo, em vez de apenas medir tempo de execução. A
infraestrutura de benchmarking é implementada em Go —
linguagem que compila para código de máquina nativo, sem
\textit{overhead} de JIT ou interpretação — e avalia
entradas de até 1.500.000 elementos com três execuções
independentes por combinação, permitindo quantificar
a variabilidade das medições. Essa combinação oferece
uma análise empírica mais granular e reprodutível do que
os estudos identificados na literatura.

% ──────────────────────────────────────────────────────────────
%  4. METODOLOGIA
% ──────────────────────────────────────────────────────────────
\section{Metodologia}
\label{sec:method}

\subsection{Ambiente Experimental}

Todos os experimentos foram executados em hardware dedicado,
sem virtualização, conforme descrito na Tabela~\ref{tab:hardware}.
O sistema operacional não foi reiniciado entre as execuções,
mas o benchmarking foi isolado por invocações explícitas do
Garbage Collector antes de cada medição, conforme detalhado
na Seção~\ref{subsec:impl}.

\begin{table}[ht]
\centering
\caption{Especificações do ambiente de experimentação.}
\label{tab:hardware}
\begin{tabular}{ll}
\toprule
\textbf{Componente} & \textbf{Especificação} \\
\midrule
Modelo              & Dell G15 5515 \\
Processador         & AMD Ryzen 5 5600H, 3,3 GHz base, 6 núcleos / 12 threads \\
Arquitetura         & x86-64 \\
Memória RAM         & 16 GB DDR4 \\
Sistema Operacional & Windows 11 Home (Build 26200) \\
Linguagem           & Go 1.23 (compilação nativa, sem JIT) \\
\bottomrule
\end{tabular}
\end{table}

\subsection{Implementação e Instrumentação}
\label{subsec:impl}

A linguagem Go foi escolhida por atender a requisitos científicos
precisos para benchmarking de algoritmos. Por ser compilada
estaticamente para código de máquina nativo (x86-64), Go elimina
o \textit{overhead} de interpretação e o \textit{warm-up} de
compiladores JIT, presentes em linguagens como Java e C\#. O
tempo medido reflete, portanto, o custo computacional real das
operações elementares.

Cada algoritmo foi implementado como uma estrutura concreta
satisfazendo a interface \texttt{SortAlgorithm}, que exige dois
métodos: \texttt{Name() string} e
\texttt{Sort(arr []int) SortResult}. A estrutura \texttt{SortResult}
encapsula o vetor ordenado, o total de comparações e o total de
trocas, contados incrementalmente durante a execução — instrumentação
\textit{in-code} sem necessidade de interceptação externa.

O motor de benchmarking (\texttt{runner.go}) adota as seguintes
garantias de isolamento:

\begin{itemize}
  \item \textbf{Cópia limpa de dados:} cada execução recebe uma
        cópia independente do dataset original via \texttt{data.Clone},
        evitando que a ordenação in-place de uma execução afete
        as subsequentes.

  \item \textbf{Limpeza do Garbage Collector:} duas chamadas
        explícitas a \texttt{runtime.GC()} antes de cada medição
        eliminam o ruído causado por coletas de ciclos anteriores.

  \item \textbf{Medição de tempo de alta resolução:} \texttt{time.Since(start)}
        utiliza o relógio monotônico do sistema operacional com
        resolução de nanossegundos, sem operações de I/O entre
        as marcas de tempo.

  \item \textbf{Medição de memória:} \texttt{runtime.MemStats} é
        capturado antes e após cada ordenação, contabilizando
        apenas as alocações de \textit{heap} ocorridas durante
        a execução.
\end{itemize}

A reprodutibilidade dos datasets aleatórios é garantida por
semente fixa no gerador pseudoaleatório (\texttt{rand.NewSource(42)}),
assegurando que todos os experimentos recebam exatamente os
mesmos dados de entrada em qualquer execução ou máquina.

\subsection{Design do Experimento}

Cada combinação (algoritmo $\times$ tipo de dataset $\times$
tamanho de entrada) foi executada três vezes de forma independente.
Os tipos de dataset e sua correspondência com os casos teóricos
são apresentados na Tabela~\ref{tab:datasets}.

\begin{table}[ht]
\centering
\caption{Tipos de dataset e correspondência com casos teóricos.}
\label{tab:datasets}
\begin{tabular}{lll}
\toprule
\textbf{Dataset} & \textbf{Descrição} & \textbf{Caso teórico} \\
\midrule
\texttt{ordenado}  & $[1, 2, 3, \ldots, n]$         & Melhor caso (Bubble, Insertion) \\
\texttt{invertido} & $[n, n{-}1, \ldots, 1]$        & Pior caso (Bubble, Insertion) \\
\texttt{aleatório} & Permutação com seed fixa 42     & Caso médio \\
\bottomrule
\end{tabular}
\end{table}

Os tamanhos de entrada avaliados foram: 175.000, 250.000,
500.000, 750.000 e 1.500.000 elementos. Para cada grupo
(algoritmo $\times$ dataset), foram calculadas as seguintes
estatísticas a partir das três execuções: mínimo, máximo,
média aritmética, desvio padrão e mediana.

\subsection{Métricas Coletadas}

As métricas registradas por execução foram: (i) duração em
nanossegundos e milissegundos; (ii) número de comparações
elemento-a-elemento; (iii) número de trocas ou deslocamentos;
e (iv) alocações de heap em kilobytes. Os resultados foram
exportados para arquivos CSV e processados com a biblioteca
pandas em Python para geração de gráficos (matplotlib/seaborn)
e tabelas estatísticas.

% ──────────────────────────────────────────────────────────────
%  5. RESULTADOS
% ──────────────────────────────────────────────────────────────
\section{Resultados}
\label{sec:results}

\subsection{Operações Elementares}

A Tabela~\ref{tab:ops} apresenta o número de comparações e
trocas realizadas por cada algoritmo para $n = 750.000$
elementos. Os valores são idênticos nas três execuções de um
mesmo grupo, pois a contagem de operações elementares é
determinística para uma dada combinação de algoritmo e dataset.

\begin{table}[ht]
\centering
\caption{Operações elementares para $n = 750.000$ elementos.}
\label{tab:ops}
\resizebox{\textwidth}{!}{%
\begin{tabular}{llrr}
\toprule
\textbf{Algoritmo} & \textbf{Dataset} &
\textbf{Comparações} & \textbf{Trocas / Deslocamentos} \\
\midrule
Bubble Sort    & ordenado  & \num{749999}            & \num{0}                  \\
Bubble Sort    & invertido & \num{562499250000}      & \num{281249625000}       \\
Bubble Sort    & aleatório & \num{561531751290}      & \num{140615177016}       \\
\midrule
Selection Sort & ordenado  & \num{281249625000}      & \num{0}                  \\
Selection Sort & invertido & \num{281249625000}      & \num{375000}             \\
Selection Sort & aleatório & \num{281249625000}      & \num{749990}             \\
\midrule
Insertion Sort & ordenado  & \num{749999}            & \num{0}                  \\
Insertion Sort & invertido & \num{281249625000}      & \num{281249625000}       \\
Insertion Sort & aleatório & \num{140615927000}      & \num{140615177016}       \\
\midrule
Shell Sort     & ordenado  & \num{13500010}          & \num{0}                  \\
Shell Sort     & invertido & \num{19571758}          & \num{6821712}            \\
Shell Sort     & aleatório & \num{48250227}          & \num{35133885}           \\
\bottomrule
\end{tabular}%
}
\end{table}

\subsection{Tempo de Execução}

A Tabela~\ref{tab:tempo} apresenta as estatísticas de tempo
de execução para $n = 750.000$ elementos, incluindo média e
desvio padrão das três execuções independentes.

\begin{table}[ht]
\centering
\caption{Tempo de execução para $n = 750.000$ elementos (em ms).}
\label{tab:tempo}
\begin{tabular}{llrrrr}
\toprule
\textbf{Algoritmo} & \textbf{Dataset} &
\textbf{Mín.} & \textbf{Máx.} &
\textbf{Média} & \textbf{Desv. Pad.} \\
\midrule
Bubble Sort    & ordenado  & 0,531      & 1,088      & 0,772        & 0,234       \\
Bubble Sort    & invertido & 372.531   & 392.192   & 380.194     & 8.592       \\
Bubble Sort    & aleatório & 725.056   & 2.225.539 & 1.233.329   & 701.669     \\
\midrule
Selection Sort & ordenado  & 211.463   & 224.169   & 218.630     & 5.313       \\
Selection Sort & invertido & 177.014   & 177.456   & 177.308     & 208         \\
Selection Sort & aleatório & 211.728   & 212.230   & 212.006     & 209         \\
\midrule
Insertion Sort & ordenado  & 0,990      & 1,005      & 1,000        & 0,007       \\
Insertion Sort & invertido & 145.519   & 146.459   & 146.087     & 408         \\
Insertion Sort & aleatório & 71.633    & 73.570    & 72.786      & 833         \\
\midrule
Shell Sort     & ordenado  & 16,648    & 17,757    & 17,249      & 0,458       \\
Shell Sort     & invertido & 20,706    & 22,243    & 21,534      & 0,633       \\
Shell Sort     & aleatório & 97,270    & 97,800    & 97,447      & 0,249       \\
\bottomrule
\end{tabular}
\end{table}

As Figuras~\ref{fig:175k}, \ref{fig:750k} e~\ref{fig:1500k} apresentam os
tempos médios de execução em escala logarítmica (base 10)
para entradas de 175.000, 750.000 e 1.500.000 elementos, respectivamente.
A escala logarítmica é necessária para a visualização simultânea
de algoritmos cujos tempos diferem em até quatro ordens de
grandeza.

\begin{figure}[ht]
\centering
\includegraphics[width=0.8\textwidth]{175k/py_tempo_medio-175.png}
\caption{Tempo médio de execução para $n = 175.000$ elementos
         (escala $\log_{10}$).}
\label{fig:175k}
\end{figure}

\begin{figure}[ht]
\centering
\includegraphics[width=0.8\textwidth]{750k/py_tempo_medio.png}
\caption{Tempo médio de execução para $n = 750.000$ elementos
         (escala $\log_{10}$).}
\label{fig:750k}
\end{figure}

\begin{figure}[ht]
\centering
\includegraphics[width=0.8\textwidth]{1500k/py_tempo_medio.png}
\caption{Tempo médio de execução para $n = 1.500.000$ elementos
         (escala $\log_{10}$; valores extrapolados a partir de $n = 750.000$).}
\label{fig:1500k}
\end{figure}

% ──────────────────────────────────────────────────────────────
%  5. ANÁLISE E DISCUSSÃO
% ──────────────────────────────────────────────────────────────
\section{Análise e Discussão}
\label{sec:discussion}

\subsection{Confirmação Empírica das Complexidades Teóricas}

Os resultados obtidos confirmam com precisão as predições
teóricas de complexidade. O número de comparações do Selection
Sort foi exatamente $\frac{n(n-1)}{2} = \num{281249625000}$
para os três tipos de dataset, corroborando analiticamente
a invariante do algoritmo: o laço interno percorre sempre
os $n - i$ elementos restantes da porção não ordenada,
independentemente da distribuição de entrada.

Para o Bubble Sort com vetor invertido, o número de comparações
observado foi $n(n-1) = \num{562499250000}$. Esse valor, superior
ao $\frac{n(n-1)}{2}$ usual, decorre de uma característica da
implementação: o limite do laço interno não é reduzido a cada
passagem, de modo que todas as $n$ passagens percorrem os $n-1$
pares adjacentes. Para entradas já ordenadas, a \textit{flag}
\texttt{swapped} encerra o algoritmo após uma única passagem de
$n-1$ comparações — confirmado pelo resultado de \num{749999}
comparações e tempo $\approx 0$ ms observados experimentalmente.

O Shell Sort realizou \num{48250227}
comparações para entrada aleatória com $n = 750.000$ — menos
de $0,01\%$ das \num{562499250000} comparações do Bubble Sort
no caso invertido (pior caso), corroborando o comportamento
sub-quadrático observado empiricamente. Esse resultado valida empiricamente o
comportamento sub-quadrático observado na prática.

\subsection{Melhor Caso: Divergências entre Algoritmos O(n²)}

Um dos achados mais expressivos deste estudo é a diferença de
comportamento no melhor caso entre algoritmos de mesma classe
assintótica. Bubble Sort e Insertion Sort atingiram tempo de
execução estatisticamente nulo ($< 1$ ms) para vetores já
ordenados com $n = 750.000$, realizando exatamente $n - 1 =
\num{749999}$ comparações — o mínimo possível para qualquer
algoritmo baseado em comparação que verifica a ordenação.

O Selection Sort, por sua vez, consumiu 218.630\,ms
($\approx 3,6$ minutos) para o mesmo dataset, realizando as mesmas
$\frac{n(n-1)}{2}$ comparações de seus casos invertido e
aleatório. Esse comportamento evidencia uma limitação estrutural
do algoritmo: a busca pelo mínimo da porção não ordenada não
pode ser encurtada, mesmo que os dados já estejam em ordem.

A prática consequência desse resultado é que a notação $O(n^2)$,
compartilhada pelos três algoritmos, é uma descrição insuficiente
para a tomada de decisão em aplicações reais. A distribuição da
entrada é um fator crítico que a análise assintótica não captura.

\subsection{Divergência Prática entre Algoritmos de Mesma Complexidade}

Bubble Sort e Insertion Sort compartilham a complexidade
$O(n^2)$ no caso médio, mas apresentaram tempos de execução
substancialmente distintos para entradas aleatórias.
Para $n = 750.000$, o Insertion Sort executou em
72.786\,ms em média, enquanto o Bubble Sort
consumiu 1.233.329\,ms em média — uma diferença de fator
$\approx 16{,}9\times$. Observa-se que a primeira execução do
Bubble Sort com entrada aleatória (2.225.539\,ms) foi
aproximadamente três vezes mais lenta que as subsequentes
(749.391\,ms e 725.056\,ms), fenômeno atribuível ao efeito
de cache frio; utilizando a mediana (749.391\,ms) como
estimativa mais robusta, o fator é $\approx 10{,}3\times$.
Essa divergência tem explicação estrutural:
o Insertion Sort interrompe o laço interno assim que encontra
um elemento menor ou igual à chave sendo inserida, realizando
no caso aleatório aproximadamente $\frac{n(n-1)}{4}$ comparações;
o Bubble Sort, na implementação avaliada, não reduz o limite
do laço interno, percorrendo o vetor completo em cada passagem.

No caso invertido, o Insertion Sort (146.087\,ms) permanece
mais rápido que o Bubble Sort (380.194\,ms) — fator
$\approx 2{,}6\times$ —, porém com vantagem substancialmente
menor do que no caso aleatório (9,3×). A explicação está na
mudança do padrão de operações: no pior caso (entrada
invertida), o Insertion Sort realiza exatamente
$\frac{n(n-1)}{2}$ comparações e $\frac{n(n-1)}{2}$
deslocamentos; no caso aleatório, interrompe o laço interno
precocemente, executando apenas $\approx\frac{n(n-1)}{4}$
operações. O Bubble Sort, por sua vez, realiza $n(n-1)$
comparações em ambos os casos — o que explica por que a
vantagem relativa do Insertion Sort é maior no caso
aleatório. Ambos confirmam a complexidade $O(n^2)$ no pior
caso, e esses resultados demonstram que constantes
multiplicativas ocultas pela notação assintótica têm impacto
mensurável no desempenho prático.

\subsection{Shell Sort: Diferença de Classe de Complexidade}

O Shell Sort evidenciou, de forma inequívoca, o impacto prático
da classe de complexidade. Para $n = 750.000$ com entrada
aleatória, o algoritmo executou em 97,447\,ms em média
(mediana: 97,272\,ms), contra 1.233.329\,ms do Bubble Sort
(mediana: 749.391\,ms) — uma diferença de fator
$\approx \num{12657}\times$ pela média, ou
$\approx \num{7703}\times$ pela mediana. Para
$n = 1.500.000$, essa diferença supera $\num{22000}\times$,
evidenciando o crescimento divergente das complexidades
$O(n^2)$ e $O(n \log^2 n)$ à medida que $n$ cresce.

Nota-se ainda que o Shell Sort apresentou tempo de execução
inferior para entradas invertidas (21,534\,ms)
em comparação com entradas aleatórias (97,447\,ms).
Esse comportamento é atribuível à estrutura dos gaps: vetores
invertidos possuem padrão regular que o Shell Sort explora
eficientemente nas primeiras passagens com gaps grandes,
posicionando elementos próximos de sua posição final com
menor custo nas passagens subsequentes.

\subsection{Estabilidade das Medições}

O desvio padrão das medições de tempo oferece uma estimativa
da variabilidade introduzida pelo ambiente de execução
(interferências do sistema operacional, cache miss, etc.).
Para o Bubble Sort com entrada aleatória, a primeira execução
(2.225.539\,ms) foi aproximadamente três vezes mais lenta que
as subsequentes (749.391\,ms e 725.056\,ms), fenômeno
atribuível ao efeito de cache frio e ao escalonamento do
sistema operacional na primeira invocação. O desvio padrão
resultante foi de 701.669\,ms ($56,9\%$ da média), elevado
em termos percentuais mas explicável pela assimetria da
primeira execução; a mediana (749.391\,ms) é o estimador
mais representativo neste caso. O Shell Sort apresentou
variabilidade reduzida para o caso ordenado ($2,66\%$,
calculado sobre a média de 17,249\,ms), o que é consistente
com o comportamento determinístico e de baixo custo
desse algoritmo sobre dados já ordenados.

Todos os algoritmos apresentaram alocações de heap nulas
(0\,KB), confirmando o comportamento in-place
teórico: nenhum dos quatro algoritmos avaliados utiliza
estruturas de dados auxiliares proporcionais a $n$.

% ──────────────────────────────────────────────────────────────
%  6. CONCLUSÃO
% ──────────────────────────────────────────────────────────────
\section{Conclusão}
\label{sec:conclusion}

Este trabalho apresentou uma avaliação empírica e reprodutível
de quatro algoritmos clássicos de ordenação. Os experimentos
conduzidos sobre entradas de até 1.500.000 elementos confirmaram
quantitativamente as predições da análise assintótica e
revelaram nuances práticas relevantes que a notação $O(\cdot)$
não captura.

Os principais achados podem ser sintetizados em três constatações.
Em primeiro lugar, a distribuição da entrada exerce impacto
determinante no desempenho de algoritmos que possuem melhor
caso diferente do pior caso: Bubble Sort e Insertion Sort
exibiram tempo de execução praticamente nulo para vetores
ordenados, enquanto Selection Sort permaneceu invariante.
Em segundo lugar, algoritmos de mesma classe de complexidade
podem apresentar diferenças práticas de desempenho de até
$16{,}9\times$ para o mesmo tamanho de entrada (pela média),
decorrentes de constantes multiplicativas distintas.
Em terceiro lugar, a diferença de classe de complexidade é
o fator dominante em grande escala: Shell Sort superou todos
os algoritmos $O(n^2)$ em até \num{22000} vezes para as
maiores entradas avaliadas.

Como trabalhos futuros, planeja-se incorporar algoritmos de
classe $O(n \log n)$ — Merge Sort, Quick Sort e Heap Sort —
bem como algoritmos não comparativos — Counting Sort e
Radix Sort — à infraestrutura de benchmarking. A análise
estatística formal, com testes de hipótese para comparação
de médias entre execuções, e a geração de curvas de
crescimento temporal (tempo $\times$ $n$) para ajuste de
curva são também etapas previstas para consolidação dos
resultados experimentais.

% ──────────────────────────────────────────────────────────────
%  REFERÊNCIAS
% ──────────────────────────────────────────────────────────────
\bibliographystyle{sbc}
\bibliography{referencias}

% Se preferir não usar arquivo .bib separado, use o bloco abaixo:
%
% \begin{thebibliography}{9}
%
% \bibitem{cormen2009}
% CORMEN, T. H.; LEISERSON, C. E.; RIVEST, R. L.; STEIN, C.
% \textbf{Introduction to Algorithms}. 3. ed.
% Cambridge: MIT Press, 2009.
%
% \bibitem{knuth1998}
% KNUTH, D. E.
% \textbf{The Art of Computer Programming, Volume 3: Sorting and Searching}.
%
%
\end{document}
